\chapter{\MakeUppercase{Introduction}}

\section{Motivation Behind Numerical Methods}

There are two basic approaches one may apply while trying to solve any mathematical problem: solving analytically or numerically. The process of solving analytically refers to calculating the exact solution of the given problem. It is possible to calculate the solution of many classical mathematical problems analytically and thus to derive the exact solution.

Nevertheless, there are many problems for which the analytical solution either does not exist at all or its calculation is excessively expensive. For instance, nonlinear and systems of nonlinear equations often cannot be solved analytically; thus, other methods have to be used.

Unlike the analytic approach, numeric methods allow deriving an approximate solution with a required precision. These methods utilize iterative algorithms that make it possible to gradually improve an approximation up until the required precision is obtained.

Different numerical methods have their own advantages and disadvantages. These advantages and disadvantages can be their convergence rates, computational costs and stability. So, selecting an appropriate method for a given problem has an important role. \cite{burden, atkinson}.

\section{Nonlinear Equations and Systems}

A nonlinear equation is an equation of the form
\begin{equation}
f(x) = 0,
\end{equation}
\noindent where \( f \) is a nonlinear function. In contrast to linear equations, nonlinear equations may exhibit multiple roots, no real roots, or complex solution structures, making their analysis and solution more challenging \cite{burden, atkinson}.

A simple example of a nonlinear equation is

\begin{equation}
f(x) = x^3 - x - 2 = 0
\end{equation}

for which no closed-form analytical solution exists in a simple form. Such equations commonly arise in various scientific and engineering applications, including physics, optimization, and computational modeling.

However, in practice, there exist cases where obtaining an exact analytical solution for a nonlinear equation is not feasible. If at all an analytical solution is obtainable, it might prove to be extremely complex to compute. This is why numerical solutions are sought after in order to obtain the roots of the nonlinear equation.

The problem can be extended to systems of nonlinear equations. 
A nonlinear system can be written in vector form as
\begin{equation}
F(\mathbf{x}) = \mathbf{0},
\end{equation}
where
\begin{equation}
F : \mathbb{R}^n \to \mathbb{R}^n
\end{equation}
is a vector-valued nonlinear function.

For example, consider the system
\begin{equation}
\begin{cases}
x^2 + y^2 - 4 = 0, \\
x - y - 1 = 0.
\end{cases}
\end{equation}

\noindent Such systems arise naturally in applications where multiple variables interact in a nonlinear manner.

Solving a nonlinear system is usually more complex than solving a single nonlinear equation. Moreover, solution methods often require use of Jacobian matrix. It is defined as

\begin{equation}
J_F(\mathbf{x}) =
\begin{bmatrix}
\frac{\partial f_1}{\partial x_1} & \frac{\partial f_1}{\partial x_2} & \cdots & \frac{\partial f_1}{\partial x_n} \\
\frac{\partial f_2}{\partial x_1} & \frac{\partial f_2}{\partial x_2} & \cdots & \frac{\partial f_2}{\partial x_n} \\
\vdots & \vdots & \ddots & \vdots \\
\frac{\partial f_n}{\partial x_1} & \frac{\partial f_n}{\partial x_2} & \cdots & \frac{\partial f_n}{\partial x_n}
\end{bmatrix}.
\end{equation}

which generalizes the concept of the derivative to higher dimensions \cite{kelley}.
In this thesis, both one-dimensional and multidimensional nonlinear equations will be studied. It will be particularly focused on the iterative techniques used to solve these equations numerically, taking into consideration issues such as convergence and cost.

\section{Objectives of the Thesis}

The primary aim of this dissertation is the analysis of numerical methods used in finding solutions of nonlinear equations and nonlinear equation systems. This kind of problem often occurs in different branches of science and technology; in many cases, there are no analytical solutions to such problems. For this reason, numerical methods take the lead in solving them. \cite{burden, atkinson}.

In this context, the thesis focuses on two main classes of problems. The first involves numerical methods for solving single nonlinear equations of the form \( f(x) = 0 \). Classical approaches, including the bisection, secant, and Newton-type methods, are examined with emphasis on their basic principles and convergence characteristics.

The second part of the study considers systems of nonlinear equations of the form \( F(\mathbf{x}) = \mathbf{0} \). Compared with single equations, these systems are generally more complex and often require more advanced solution techniques, especially those based on the Jacobian matrix. The methods used to solve such systems are investigated within a similar theoretical framework.

In addition to the theoretical analysis, the selected numerical methods are implemented in Python using a consistent computational structure. The implementations are used to investigate the practical behavior of the methods through numerical experiments. In particular, the methods are compared in terms of convergence behavior, residual reduction, and robustness.

Overall, the aim of this thesis is to provide a clear and structured examination of numerical methods for nonlinear problems, combining theoretical foundations with computational analysis.

\section{Structure of the Thesis}

This thesis is divided into eight chapters.

Chapter 1 explains the motivation for studying numerical methods and introduces the concepts of nonlinear equations and systems of nonlinear equations. The aims of the thesis and the topics covered throughout the study are also outlined.

Chapter 2 provides the mathematical background needed for the methods discussed in the following chapters. Concepts such as fixed-point formulations, convergence analysis, orders of convergence, norms, Jacobian matrices, and stopping criteria are presented as the foundation for the later discussions.

Chapter 3 deals with numerical methods for solving single nonlinear equations. Well-known methods, including the bisection method, secant method, Newton’s method, damped Newton method, and Brent’s method, are introduced and discussed in terms of their convergence behavior. The chapter also includes illustrative examples and examples of situations in which these methods may fail.

Chapter 4 turns to systems of nonlinear equations. Multivariable Newton’s method and Broyden’s method are examined from both theoretical and computational perspectives. Several examples are provided to illustrate their application and to compare their convergence performance.

Chapter 5 focuses on the practical implementation of the numerical methods considered in the thesis. The structure of the repository, the Python implementations, the testing procedures, and issues related to reproducibility are discussed.

Chapter 6 presents the numerical experiments and comparative analyses carried out in this study. Benchmark problems, convergence plots, residual analyses, and experimental results are reported for both single-variable equations and systems of nonlinear equations, allowing the performance of the methods to be evaluated under different conditions.

Chapter 7 explores the connection between nonlinear equation solving and optimization. Topics such as residual minimization, Newton-type optimization methods, damping strategies, line-search techniques, and the Armijo condition are discussed from an optimization viewpoint. The chapter also includes a numerical investigation of an adaptive damped Newton method that employs Armijo backtracking.

Chapter 8 concludes the thesis by summarizing the main results obtained. The limitations of the work are discussed and several possible directions for future research are suggested.